\documentclass[12pt]{article}
\usepackage{amsmath}
\usepackage{amssymb}
\usepackage{amsthm}
\usepackage{geometry}
\geometry{margin=1in}
\usepackage{xcolor}
\usepackage{asymptote}
\usepackage{lastpage}
\usepackage{fancyhdr}
\pagestyle{fancy}
\fancyhf{}
\renewcommand{\headrulewidth}{0pt}
\cfoot{\thepage\ of \pageref{LastPage}}
\fancypagestyle{plain}{%
  \fancyhf{}%
  \renewcommand{\headrulewidth}{0pt}%
  \cfoot{\thepage\ of \pageref{LastPage}}%
}

\title{AMC 12 Prep 2026}
\author{}
\date{\today}

\setcounter{secnumdepth}{0}

\begin{document}
\maketitle

\section{2008 AMC12B Problem 10}

\textbf{Problem:}
Bricklayer Brenda would take $9$ hours to build a chimney alone, and bricklayer Brandon would take $10$ hours to build it alone. When they work together they talk a lot, and their combined output is decreased by $10$ bricks per hour. Working together, they build the chimney in $5$ hours. How many bricks are in the chimney?

$\textbf{(A)}\ 500 \qquad \textbf{(B)}\ 900 \qquad \textbf{(C)}\ 950 \qquad \textbf{(D)}\ 1000 \qquad \textbf{(E)}\ 1900$

\vspace{0.3cm}
\noindent\textbf{Answer:}~\underline{\hspace{5cm}}\hfill\textbf{Time:}~\underline{\hspace{1.2cm}}:\underline{\hspace{1.2cm}}

\vspace{0.5cm}

\section{2009 AMC12A Problem 10}

\textbf{Problem:}
In quadrilateral $ABCD$ , $AB = 5$ , $BC = 17$ , $CD = 5$ , $DA = 9$ , and $BD$ is an integer. What is $BD$ ? \begin{asy}
unitsize(4mm);
defaultpen(linewidth(.8pt)+fontsize(8pt));
dotfactor=4;
pair C=(0,0), B=(17,0);
pair D=intersectionpoints(circle(C,5),circle(B,13))[0];
pair A=intersectionpoints(circle(D,9),circle(B,5))[0];
pair[] dotted={A,B,C,D};
draw(D--A--B--C--D--B);
dot(dotted);
label("$D$",D,NW);
label("$C$",C,W);
label("$B$",B,E);
label("$A$",A,NE);
\end{asy}

$\textbf{(A)}\ 11 \qquad \textbf{(B)}\ 12 \qquad \textbf{(C)}\ 13 \qquad \textbf{(D)}\ 14 \qquad \textbf{(E)}\ 15$

\vspace{0.3cm}
\noindent\textbf{Answer:}~\underline{\hspace{5cm}}\hfill\textbf{Time:}~\underline{\hspace{1.2cm}}:\underline{\hspace{1.2cm}}

\vspace{0.5cm}

\section{2009 AMC12B Problem 10}

\textbf{Problem:}
A particular $12$ -hour digital clock displays the hour and minute of a day. Unfortunately, whenever it is supposed to display a $1$ , it mistakenly displays a $9$ . For example, when it is 1:16 PM the clock incorrectly shows 9:96 PM. What fraction of the day will the clock show the correct time?

$\mathrm{(A)}\ \frac 12\qquad \mathrm{(B)}\ \frac 58\qquad \mathrm{(C)}\ \frac 34\qquad \mathrm{(D)}\ \frac 56\qquad \mathrm{(E)}\ \frac {9}{10}$

\vspace{0.3cm}
\noindent\textbf{Answer:}~\underline{\hspace{5cm}}\hfill\textbf{Time:}~\underline{\hspace{1.2cm}}:\underline{\hspace{1.2cm}}

\vspace{0.5cm}

\section{2010 AMC12A Problem 10}

\textbf{Problem:}
The first four terms of an arithmetic sequence are $p$ , $9$ , $3p-q$ , and $3p+q$ . What is the $2010^\text{th}$ term of this sequence?

$\textbf{(A)}\ 8041 \qquad \textbf{(B)}\ 8043 \qquad \textbf{(C)}\ 8045 \qquad \textbf{(D)}\ 8047 \qquad \textbf{(E)}\ 8049$

\vspace{0.3cm}
\noindent\textbf{Answer:}~\underline{\hspace{5cm}}\hfill\textbf{Time:}~\underline{\hspace{1.2cm}}:\underline{\hspace{1.2cm}}

\vspace{0.5cm}

\section{2010 AMC12B Problem 10}

\textbf{Problem:}
The average of the numbers $1, 2, 3,\cdots, 98, 99,$ and $x$ is $100x$ . What is $x$ ?

$\textbf{(A)}\ \dfrac{49}{101} \qquad \textbf{(B)}\ \dfrac{50}{101} \qquad \textbf{(C)}\ \dfrac{1}{2} \qquad \textbf{(D)}\ \dfrac{51}{101} \qquad \textbf{(E)}\ \dfrac{50}{99}$

\vspace{0.3cm}
\noindent\textbf{Answer:}~\underline{\hspace{5cm}}\hfill\textbf{Time:}~\underline{\hspace{1.2cm}}:\underline{\hspace{1.2cm}}

\vspace{0.5cm}

\end{document}